\section{Background: one pod, two ways to submit it}
\label{sec:background}

\paragraph{Requests, limits and cgroups.}
Each Kubernetes container declares a memory and CPU \emph{request}, which the scheduler
reserves on a node, and a \emph{limit}, which the kubelet enforces through the
container's cgroup~\citep{k8sresources,cgroupv2}. Under cgroup~v2, a container whose
charged memory reaches its limit is OOM-killed. Both anonymous memory and shared memory
(\code{shmem}) count toward that limit, while clean page cache does not, because the
kernel can reclaim it first. A CPU limit works differently. It is enforced as a CFS
quota, so a container that exceeds it is slowed down rather than killed. Accelerators are
requested as whole devices on node pools that carry them. In this paper a request also
sets the limit, which is the usual choice for batch pods because it makes scheduling and
failures predictable.

\paragraph{How the pod spec changes a task's memory.}
Three details of a pod spec can change how much memory a batch task needs without any
change to its code. The first is the medium of the volume that holds its inputs. An
\code{emptyDir} on node disk keeps files as page cache, which the kernel reclaims under
pressure. An \code{emptyDir} with \code{medium: Memory} is a tmpfs, and every byte written
to it is charged to memory as \code{shmem}~\citep{k8svolumes}. Staging a 2\,GiB input on
tmpfs therefore adds 2\,GiB to the pod's need. The second detail is who fetches the
inputs. An init container runs to completion before the main container starts and has
its own limit, whereas a download inside the job runs in the job's container, so the
client library's memory and threads count against the job. The third is anything else
that shares the container, such as an orchestration runtime or a client library.

There is one subtlety with tmpfs. A page is charged to the cgroup of the container that
writes it, so pages written by an init container are not charged to the main container.
When the init container exits, though, its pages move to the pod's cgroup and still count
against the pod's overall limit. We confirm this in real pods in
\cref{sec:validation}.

\paragraph{Raw Jobs.}
A \code{batch/v1} Job runs a pod template until it succeeds. Its \code{backoffLimit}
says how many failed pods it will replace before giving up. Batch teams often set it to
zero and resubmit failures by hand or from an outer scheduler. A Job has to stage its own
inputs, either with an init container or in the job's code.

\paragraph{Flyte as an interface to Kubernetes.}
Flyte~\citep{flyte} turns a typed Python function into a Kubernetes pod. A user writes
\code{@env.task async def f(x: File) -> dict}, and Flyte builds the image, renders the
pod spec, moves typed inputs and outputs through object storage, caches results by their
inputs, retries failed attempts, and records each execution so that a long pipeline can
resume after a node is lost. None of this changes what the function does inside the
container. It changes the pod's edges instead. Inputs arrive through Flyte's own
download, the container also runs Flyte's runtime, and the execution history reports the
container's peak with that runtime included. We measure the runtime at 55\,MiB at rest,
and a median 70\,MiB above the same workload run as a plain Job (\cref{sec:validation}).

For resource sizing, then, Flyte is a convenience layer over the same pod. It costs a
little memory in exchange for typed I/O, caching, retries and durable execution. An
approach that can size a Kubernetes pod should be able to size a Flyte task using the
same knowledge plus the runtime.
